\documentclass[12pt,a4paper]{article}
\usepackage{amsmath,amssymb,amsthm}
\usepackage{booktabs}
\usepackage{hyperref}
\usepackage{geometry}
\geometry{margin=2.5cm}

\newtheorem{theorem}{Theorem}[section]
\newtheorem{lemma}[theorem]{Lemma}
\newtheorem{corollary}[theorem]{Corollary}
\newtheorem{proposition}[theorem]{Proposition}
\theoremstyle{definition}
\newtheorem{definition}[theorem]{Definition}
\theoremstyle{remark}
\newtheorem{remark}[theorem]{Remark}

\title{Complex Hilbert Space from Recognition Science:\\
$e^{i\pi} = -1$ from the 8-Tick Half-Period}

\author{Jonathan Washburn\\
Recognition Science Research Institute\\
Austin, Texas\\
\texttt{washburn.jonathan@gmail.com}}

\date{March 2026}

\begin{document}
\maketitle

\begin{abstract}
We derive the necessity of complex Hilbert space for quantum
mechanics from the Recognition Science 8-tick epoch.  The epoch has
$2^D = 8$ ticks ($D = 3$).  The $k$-th tick carries phase
$e^{ik\pi/4}$ ($k = 0, 1, \ldots, 7$), the 8th roots of unity.
At the half-period ($k = 4$): $e^{i\pi} = -1$, a real quantity
acquires a sign flip.  To represent this phase evolution continuously,
complex numbers are required: $\mathbb{R}$ cannot support continuous
rotations in a plane, but $\mathbb{C}$ can.  The full-period return
$e^{2i\pi} = 1$ restores reality.  Thus $\mathbb{C}$ is the minimum
field extension needed for the 8-tick dynamics.  Real quantum mechanics
($\mathbb{R}$) fails because it cannot represent the half-period
phase.  Quaternionic quantum mechanics ($\mathbb{H}$) is unnecessary
because $\mathbb{C}$ suffices---the 8-tick structure provides no
additional cyclic structure requiring a third imaginary unit.  The sum
of all 8 phases vanishes: $\sum_{k=0}^7 e^{ik\pi/4} = 0$ (vacuum
fluctuation cancellation).
\end{abstract}

\section{Introduction}

Why does quantum mechanics use complex numbers~\cite{Stueckelberg}?
This is one of the deepest structural questions in physics.  Adler~\cite{Adler} explored quaternionic QM; recent experiments have ruled out
real QM~\cite{Renou}.  RS~\cite{RS_Axioms} derives the answer: the 8-tick
epoch forces $\mathbb{C}$.

\section{Recognition Science Framework}
\label{sec:framework}

Recognition Science derives all physics from a single primitive,
the \emph{recognition cost functional} $J(x)$, uniquely determined
by three axioms:

\begin{definition}[RS Cost Functional]
For $x > 0$: $J(x) = \tfrac{1}{2}(x + x^{-1}) - 1$.
\end{definition}

\noindent\textbf{Axiom A1 (Normalization):} $J(1) = 0$.\\
\textbf{Axiom A2 (Recognition Composition Law, RCL):}
$J(xy) + J(x/y) = 2J(x)J(y) + 2J(x) + 2J(y)$ for all $x, y > 0$.\\
\textbf{Axiom A3 (Calibration):} $J''_{\log}(0) = 1$, where
$J_{\log}(t) := J(e^t)$.

\begin{theorem}[Cost Uniqueness]
\label{thm:unique}
The unique continuous solution to \textup{(A1)--(A3)} is
$J(x) = \tfrac{1}{2}(x + x^{-1}) - 1$.
\end{theorem}

\begin{proof}[Proof sketch]
Setting $f(t) = J_{\log}(t) + 1$ and rewriting \textup{(A2)} gives the
d'Alembert equation $f(s+t) + f(s-t) = 2f(s)f(t)$.  By the Acz\'el
classification theorem~\cite{Aczel1966}, the unique continuous solution
with $f(0) = 1$ and $f''(0) = 1$ is $f(t) = \cosh t$.
\end{proof}

\noindent The 8-tick epoch ($2^D = 8$ for $D = 3$, forced by the gap-45
synchronization $\operatorname{lcm}(8,45)=360$) carries phases
$e^{ik\pi/4}$, $k = 0, \ldots, 7$---the 8th roots of unity.  These
require $\mathbb{C}$ to represent: $\mathbb{R}$ cannot encode $e^{i\pi/4}$.

\section{The 8-Tick Phase Structure}

\begin{definition}[Tick Phases]
The phase at tick $k$ ($k = 0, 1, \ldots, 7$) is
$\omega_k = e^{ik\pi/4}$, the 8th roots of unity.
\end{definition}

\begin{remark}[Periodicity]
$\omega_k^8 = e^{8ik\pi/4} = e^{2\pi ik} = 1$ for all $k \in \{0,\ldots,7\}$.
The phase returns to unity after one complete epoch.
\end{remark}

\begin{remark}[Half-Period Phase and Fermion Sign]
$\omega_4 = e^{4i\pi/4} = e^{i\pi} = -1$ (by Euler's formula).  The
RS structural identification is that fermionic statistics correspond
to the half-epoch phase: a particle whose state acquires $-1$ under
a full 8-tick cycle (2 epochs) is identified with a spin-$1/2$ fermion.
This is an identification, not a derivation of the spin-statistics
theorem (see Remark below).
\end{remark}

\begin{remark}[Boson Phase]
$\omega_0 = e^{0} = 1$.  The RS structural identification is that
bosons carry the trivial phase (no sign flip under a full cycle),
consistent with integer spin.
\end{remark}

\begin{remark}[Spin-Statistics --- Structural Identification]
The identification $\omega_4 = -1$ (fermions) and $\omega_0 = 1$
(bosons) provides an RS structural basis for the spin-statistics
connection.  The full spin-statistics theorem (Pauli, 1940) is a deep
consequence of Lorentz invariance and local field theory, not derived
here.  The RS contribution is the identification of the $-1$ sign with
the half-period phase $e^{i\pi}$ of the 8-tick epoch: particles
associated with the 4-tick half-period are identified with fermions.
A rigorous derivation of the Spin-Statistics theorem from the RS
ledger structure is an identified open problem.
\end{remark}

\section{Why $\mathbb{C}$ Is Required}

\begin{proposition}[Reals Insufficient for 8-Tick Phases]
$\mathbb{R}$ cannot represent all 8-tick phases: $\omega_1 = e^{i\pi/4}
= (1+i)/\sqrt{2} \notin \mathbb{R}$.
\end{proposition}

\begin{proof}
$\omega_1 = \cos(\pi/4) + i\sin(\pi/4) = (1 + i)/\sqrt{2}$.  This
has a nonzero imaginary part.  No real number can represent a rotation
in a plane; $\mathbb{C}$ is the minimum extension of $\mathbb{R}$
that supports multiplication-as-rotation.
\end{proof}

\begin{proposition}[Quaternions Unnecessary --- Structural Argument]
$\mathbb{H}$ provides no additional cyclic structure for the 8-tick dynamics.
\end{proposition}

\begin{proof}[Structural argument]
The 8-tick phases $\{\omega_k = e^{ik\pi/4} : k = 0,\ldots,7\}$ form
the cyclic group $\mathbb{Z}_8$, which is a subgroup of $U(1)
\subset \mathbb{C}^\times$.  Every element of $\mathbb{Z}_8$ can be
expressed using $\mathbb{C}$ alone; no third or fourth imaginary unit
($j$, $k$ in $\mathbb{H}$) is needed to generate the 8th roots of
unity.  Quaternionic QM would introduce $SU(2)$ holonomy, providing
more structure than the single $U(1)$ of the 8-tick phases.  Since the
RS 8-tick structure is a $\mathbb{Z}_8$ subgroup of $U(1)$, $\mathbb{C}$
is sufficient and $\mathbb{H}$ is redundant for this purpose.
\end{proof}

\section{The Number Field Classification}

By the Frobenius theorem~\cite{Frobenius}, the only finite-dimensional associative
division algebras over $\mathbb{R}$ are $\mathbb{R}$, $\mathbb{C}$,
and $\mathbb{H}$.  Any quantum theory with a Hilbert space must use
one of these.  $\mathbb{R}$ fails: it cannot support continuous
phase rotations; $e^{i\pi/4}$ has no analogue in $\mathbb{R}$.
$\mathbb{H}$ is unnecessary: the 8-tick structure provides one
cyclic degree of freedom; quaternions introduce $i$, $j$, $k$ with
noncommutative multiplication, but the phases live in a single
$U(1) \subset \mathbb{C}^\times$.  $\mathbb{C}$ is uniquely
selected: the minimal division algebra that extends $\mathbb{R}$,
contains the 8th roots of unity, and introduces no superfluous
structure.

\section{Experimental Evidence Against Real QM}

Renou et al.~\cite{Renou} (2021) performed a loophole-free Bell
experiment that ruled out real quantum mechanics.  In real QM the
state space is over $\mathbb{R}$ and correlations are constrained;
the observed correlations violate the bounds real QM permits.  Real
QM predicts a maximum violation strictly less than complex QM for
certain configurations; the measured violations exceeded the
real-QM bound, confirming complex amplitudes.  This is consistent
with RS: the 8-tick structure requires $\mathbb{C}$, and experiment
excludes $\mathbb{R}$.

\section{Vacuum Fluctuation Cancellation}

\begin{theorem}[Phase Sum Vanishes]
$\sum_{k=0}^7 e^{ik\pi/4} = 0$.
\end{theorem}

\begin{proof}
The sum of all $n$th roots of unity vanishes for $n \geq 2$:
$\sum_{k=0}^{n-1} e^{2\pi ik/n} = 0$.  For $n = 8$:
$\sum_{k=0}^7 e^{ik\pi/4} = 0$.
\end{proof}

This has physical significance: vacuum fluctuations at each of the 8
tick phases cancel exactly, explaining why the vacuum energy is not
divergent but finite.

\section{The Cosmological Constant Connection}

The cosmological constant problem: naive QFT predicts vacuum energy
$\sim M_{\mathrm{Pl}}^4$, while observation gives $\rho_\Lambda \sim
10^{-120} M_{\mathrm{Pl}}^4$.  In RS, $\sum_{k=0}^7 \omega_k = 0$
provides a structural contribution to this cancellation: vacuum
fluctuations associated with the 8-tick epoch phases sum to zero
exactly.  The RS claim is that the leading-order epoch contribution to
vacuum energy vanishes by this phase cancellation.

\begin{remark}[Status of the Cosmological Constant Connection]
The full resolution of the cosmological constant problem requires
explaining why $\rho_\Lambda \ll M_{\mathrm{Pl}}^4$; the phase sum
$\sum \omega_k = 0$ shows that the 8-tick epoch contribution is exactly
zero, but does not address vacuum energy contributions from QFT modes
outside the epoch structure.  The identification of ``epoch vacuum
fluctuations'' with the dominant divergent QFT contribution is a
structural hypothesis, not a derivation.  This is an identified open
problem; a companion paper on the Cosmological Constant addresses it
in detail.
\end{remark}

\section{Euler's Formula}

\begin{remark}[Euler's Formula in RS Context]
The standard identity $e^{i\theta} = \cos\theta + i\sin\theta$
(Euler's formula) describes how the 8-tick phase $e^{ik\pi/4}$ evolves
continuously between commits.  In RS, Euler's formula is not merely a
mathematical curiosity: it is the structural necessity that relates the
discrete tick-phase assignments to continuous-time quantum dynamics.
The amplitude of a state between commits is described by a complex
exponential; the phase wraps around the unit circle over the 8-tick epoch.
\end{remark}

\section{Comparison of Quantum Formulations}

Table~\ref{tab:qm-comparison} summarizes the four formulations.

\begin{table}[htbp]
\centering
\caption{Comparison of quantum formulations.}
\label{tab:qm-comparison}
{\small
\begin{tabular}{@{}lcccc@{}}
\toprule
Property & Real & Complex & Quat.\ & RS \\
\midrule
Number field & $\mathbb{R}$ & $\mathbb{C}$ & $\mathbb{H}$ & $\mathbb{C}$ \\
8-tick phases & No & Yes & Yes & Yes \\
Half-period & No & Yes & Yes & Yes \\
$\sum \omega_k$ & N/A & 0 & 0 & 0 \\
Experiment & Ruled out & OK & N/A & OK \\
Minimal & No & Yes & No & Yes \\
\bottomrule
\end{tabular}}
\end{table}

Real QM cannot represent the 8-tick phases.  Complex QM is minimal
and experimentally confirmed.  Quaternionic QM is unnecessarily
rich.  RS predicts $\mathbb{C}$ and is consistent with experiment.

\section{Predictions}

\begin{enumerate}
\item \textbf{No quaternionic effects.}  RS predicts no observables
  sensitive to $j$ or $k$; the 8-tick structure provides no such
  degrees of freedom.
\item \textbf{Vacuum energy.}  Dominant vacuum contribution from the
  epoch vanishes; residual $\Lambda$ terms arise from other
  mechanisms.
\item \textbf{Phase coherence.}  Interference at fundamental scales
  should privilege phase differences $\Delta\theta = m\pi/4$
  ($m \in \mathbb{Z}$).
\item \textbf{No real-QM revival.}  RS and Renou et al.\ exclude
  real QM; any future support would contradict both.
\end{enumerate}

\section{Conclusion}

Complex Hilbert space is forced by the 8-tick epoch: the phases
$e^{ik\pi/4}$ require $\mathbb{C}$.  Real QM is excluded
($\omega_1$ has imaginary part).  Quaternionic QM is unnecessary
(no additional cyclic structure).  Vacuum fluctuations cancel
($\sum \omega_k = 0$).

\begin{thebibliography}{9}

\bibitem{Aczel1966}
J.~Acz\'el, \emph{Lectures on Functional Equations and Their Applications},
Academic Press, New York (1966).

\bibitem{RS_Axioms}
J.~Washburn, ``The Algebra of Reality: A Recognition Science Derivation
of Physical Law,'' \emph{Axioms} \textbf{15}(2), 90 (2026).
\texttt{doi:10.3390/axioms15020090}

\bibitem{Stueckelberg}
E.~C.~G.~Stueckelberg, ``Quantum Theory in Real Hilbert Space,''
Helv.\ Phys.\ Acta \textbf{33}, 727 (1960).

\bibitem{Adler}
S.~L.~Adler, \textit{Quaternionic Quantum Mechanics and Quantum Fields}
(Oxford University Press, 1995).

\bibitem{Renou}
M.-O.~Renou \emph{et al.}, ``Quantum Theory Based on Real Numbers
Can Be Experimentally Falsified,'' Nature \textbf{600}, 625 (2021).

\bibitem{Frobenius}
F.~G.~Frobenius, ``\"Uber lineare Substitutionen und bilineare
Formen,'' J.\ Reine Angew.\ Math.\ \textbf{84}, 1 (1878).
\end{thebibliography}

\end{document}
